\addchap{Abstract}
\label{chap:abstract}

Jordan sorting recovers the coordinate order of intersections along a line
from the order in which a simple curve encounters them. Historical
linear-time algorithms solve this problem through an incremental sorting
procedure combined with specialized dynamic-list structures. This thesis
studies whether the control flow of the 1990 Simplified Jordan Sorting
algorithm can be reconstructed as an executable, auditable implementation
when those theoretical structures are replaced by ordinary Python lists.

The reconstruction maintains its own partial sorted order, upper and lower
pair families, and sibling-list ownership while executing the paper-facing
Step~1/2/3 procedure. For oracle-certified valid inputs, the paper core
recovers its result from the maintained partial order and does not consume the
oracle's sorted output. Transactional sibling-list updates, untimed
checked-state audits, deterministic replay, and fixed execution policies separate
correctness instrumentation from the pre-certified \texttt{minimal} call used
for timing. Under a pre-specified protocol, the formal experiment evaluates 60
controlled cases at five sizes from $n=32$ to $n=512$, comparing Python sort, a
complete oracle-backed reference pipeline, and the ordinary-list paper
reconstruction.

All 3600 measured rows in the formal experiment have correct output and no
recorded error, and all 60 untimed checked-state audits pass. The median
exact-case paper/reference runtime ratio decreases from $3.226$ at $n=32$ to $0.567$ at
$n=512$; it is above one through $n=128$ and below one at $n=256$ and
$n=512$. Python sort has the lowest median call time at every tested size.
These ratios are pipeline-scope comparisons: the reference measurement
includes its complete oracle-backed pipeline, whereas the paper measurement
contains only the pre-certified \texttt{minimal} sorting call, with oracle
certification and checked diagnostics outside timing. They are therefore not
like-for-like end-to-end speedups.

The contribution is an independently recovering ordinary-list reconstruction
with explicit state, safety, audit, and evidence boundaries, together with a
reproducible empirical characterization of its behavior on the controlled
sample. The implementation does not include the specialized theoretical
backend, and the experiment does not evaluate recognition. No linear-time or
asymptotic-complexity claim is made for the Python implementation, and the
reported runtimes remain specific to the tested cases, timing scopes, and
recorded execution environment.
